\begin{titlepage}
    \section*{Abstract}
        % Set the Scene
        Single-board computers such as Raspberry Pi 4 and 5, represent compact 
        and power-efficient systems intended for diverse areas of applications. 
        These platforms depend on main memory to supply fast and temporary storage 
        for data currently being processed. In \gls{dram}, each bit is maintained within 
        an individual memory cell, where a binary state is encoded by the existence 
        or absence of electrical charge in a capacitor~\cite{paperFlippingBits, paperRowhammeronaRPi, paperDrammer}.
        % Conflict
        As \gls{dram} cells are further miniaturized to achieve higher memory densities 
        within constrained physical space, advances in technology improve storage 
        capacity and energy efficiency but also increase vulnerability to disturbance 
        phenomena. Reliable \gls{dram} requires maintaining data integrity despite charge 
        leakage. Repeated gate-induced drain- and subthreshold leakage~\cite{paperFlippingBits, paperRHinDRAM} can 
        accelerate charge loss and intensify cell-to-cell interference, eventually 
        causing bits to flip from 0 to 1 or vice versa. Rowhammer leverages this 
        behavior by intentionally provoking rapid, repeated row activations that 
        intensify charge loss and disturbance errors. As a result, memory contents 
        may be corrupted, potentially facilitating privilege escalation and other 
        critical security impacts.
        % Experiments
        In this work, we develop a Rowhammer proof-of-concept for the Raspberry 
        Pi 4 and 5. the implementation enables direct \gls{dram} access by employing cache-bypass 
        techniques based on eviction strategies and ARM cache maintenance instructions.
        To facilitate targeted row accesses, \gls{dram} address mapping was reconstructed by ported 
        ARM-compatible DRAMA tool.
        % Climax Experiment
        Additionally, TRRespass was adapted to ARM to evaluate effective access patterns 
        for hammering designed to induce bitflips.
        % Validate Climax Experiment
        The proof-of-concept integrates all functional components, ranging from 
        cache-bypass mechanisms and targeted \gls{dram} access to effective hammering 
        patterns, into a framework designed to experimentally assess susceptibility 
        to Rowhammer on contemporary Raspberry Pi hardware. While the full proof-of-concept 
        successfully executed on multiple devices and RAM configurations, no 
        reproducible bitflips were observed. Hardware protections and \gls{dram} page 
        policies limit the effectiveness of hammering patterns, demonstrating the 
        practical challenges in exploiting Rowhammer on these platforms.
        \\
        \textit{\textbf{Keywords}: Rowhammer, \gls{dram}, Raspberry Pi, Address Mapping, \gls{trr}}
\end{titlepage}
